\input{../../../headers/tpheaders_2.tex}


\prob{Modéliser les caractéristiques d'un moteur à courant continu.} 

\\

\section{Modèle du moteur électrique à courant continu}

\begin{eqnarray}
u_m(t)=L_m.\frac{di(t)}{dt}+R_m.i(t)+e(t) \\
e(t)=K_e.\omega_m(t) \\
J.\frac{d\omega_m(t)}{dt}=C_m(t)-C_r(t) \\
C_m(t)=K_m.i(t)
\end{eqnarray}

Données:
\begin{itemize}
 \item $u_m(t)$: tension aux bornes du moteur ($V$),
 \item $i(t)$: intensité du courant dans le moteur ($A$),
 \item $e(t)$: force électromotrice ($V$),
 \item $\omega_m(t)$: vitesse de rotation du moteur ($rad.s^{-1}$),
 \item $C_m(t)$: couple moteur ($N.m$),
 \item $C_r(t)$: couple résistant ($N.m$),
 \item $L_m$: inductance de la bobine du moteur ($H$),
 \item $R_m$: résistance électrique interne au moteur ($\Omega$),
 \item $K_e$: constante électrique du moteur ($V.rad^{-1}.s$),
 \item $J$: inertie du moteur ($kg.m^2$),
 \item $K_m$: constante de couple du moteur ($N.m.A^{-1}$).
\end{itemize}

~\

En général, on suppose $K_e=K_m$ pour une MCC.

\paragraph{Question 1:} D'après les équations (1) à (4), \textbf{écrire} une équation liant $u_m(t)$, $\omega_m(t)$ et $C_m(t)$.

\paragraph{Question 2:} Quelles sont les \textbf{hypothèses} à poser afin de mettre cette équation sous la forme $u_m(t)=K.\omega_m(t)$.

\paragraph{Question 3:} En supposant que l'on arrive à mesurer le courant qui traverse le moteur, que devient-il alors possible d'\textbf{estimer} ?

\newpage

\section{Mesure des valeurs caractéristiques du moteur} 

\paragraph{Question 4:} \textbf{Déterminer} un protocole afin de mesurer la tension aux bornes du moteur. Donner la liste du matériel utilisé et ses caractéristiques (sensibilité, plage de mesure...).

\paragraph{Question 5:} \textbf{Déterminer} un protocole afin de mesurer la vitesse de rotation du moteur à l'aide du tachymètre. Donner les caractéristiques du matériel (sensibilité, plage de mesure...).

\paragraph{Question 6:} \textbf{Mettre en \oe uvre} ce protocole pour des tensions allant de 0V à 12V. Écrire les valeurs mesurés dans le script 
\href{https://github.com/Costadoat/Sciences-Ingenieur/raw/master/S01\%20Analyse\%20fonctionnelle/TP01\%20Modélisation\%20linéaire/Ilot_01\%20Moteur\%20\%C3\%A0\%20courant\%20continu/01-TP01-I01.py}{te} python
un tableur et tracer la courbe $\omega_m(t)=f(u_m(t))$. Conclure quant à l'allure de ce tracé.

\paragraph{Question 7:} \textbf{Refaire} la mesure précédente en mesurant aussi le courant qui traverse le moteur.

\paragraph{Question 8:} \textbf{Proposer} et \textbf{mettre en \oe uvre} un protocole permettant de mesurer la résistance interne du moteur. Donner la liste du matériel utilisé et ses caractéristiques (sensibilité, plage de mesure...).

\section{Vérification des modèles et analyse des résultats}

\paragraph{Question 9:} À partir des résultats précédents, \textbf{déterminer} le couple résistant $C_r(t)$ pour le moteur libre. \textbf{Proposer} une solution permettant d'augmenter ce couple, \textbf{prédire} le comportement du système.

\subsection{Vérifier expérimentalement un modèle théorique}

\paragraph{Question 10:} \textbf{Mettre en \oe uvre} ce nouveau protocole et \textbf{conclure} quant à la validité de la prédiction de la question 9.

\subsection{Déterminer la valeur de $L_m$}

\paragraph{Question 11:} \textbf{Proposer} sans les mettre en \oe uvre des protocoles permettant de déterminer la valeur de $L_m$.

\section{Synthèse du TP}

\paragraph{Question 12:} \textbf{Conclure} quant au modèle obtenu pour ce moteur à courant continu, \textbf{réaliser} une synthèse de ce TP présentant votre démarche pour répondre à la problématique.

\ifdef{\public}{\end{document}}{}

\clearpage

\newpage

\pagestyle{correction}

\section{Correction}

\paragraph{Question 1:} D'après les équations (1) à (4), \textbf{écrire} une équation liant $u_m(t)$, $\omega_m(t)$ et $C_m(t)$.

$u_m(t)=\frac{L_m}{K_m}.\frac{dC_m(t)}{dt}+\frac{R_m}{K_m}.C_m(t)+K_e.\omega_m(t)$

$u_m(t)=\frac{L_m.J}{K_m}.\frac{d^2\omega_m(t)}{dt^2}+\frac{R_m.J}{K_m}.\frac{d\omega_m(t)}{dt}+K_e.\omega_m(t)$


\paragraph{Question 2:} Quelles sont les \textbf{hypothèses} à prendre afin de mettre cette équation sous la forme $u_m(t)=K.\omega_m(t)$.

Il faut se placer à vitesse constante (régime établi) et négliger les frottements, ainsi, $C_r(t)=0$, $\frac{d^2\omega_m(t)}{dt^2}=0$ et $\frac{d\omega_m(t)}{dt}=0$.

Ainsi, $u_m(t)=K_e.\omega_m(t)$

\paragraph{Question 3:} En supposant que l'on arrive à mesurer le courant qui traverse le moteur, que devient-il alors possible de \textbf{mesurer} ?

Il est alors possible de déterminer le couple résistant $C_r(t)$ en régime établi.

Avec $C_m(t)=K_m.i(t)$ et $J.\frac{d\omega_m(t)}{dt}=C_m(t)-C_r(t)=0$ en régime établi, on obtient $C_r(t)=K_m.i(t)=K_e.i(t)$.

\end{document}
